\documentclass[12pt]{article}
\usepackage[utf8]{inputenc}
\usepackage[T1]{fontenc}
\usepackage{hyperref}
\title{Keiyo Language Resources for K-12 Curriculum Development}
\author{Compiled by Claude Sonnet 5.5 (claude-sonnet-5-5) for Andr\'as Kornai}
\date{}
\begin{document}
\maketitle

\section{Introduction}
Keiyo (ISO 639-3 \texttt{eyo}) is a Kalenjin-cluster Nilotic language of Kenya. This report lists only what a web search turned up. No Keiyo-specific dictionary, grammar or school material was found.

\section{Primary Educational Resources}
None found.

\section{Linguistic and Reference Materials}
\begin{itemize}
\item Wikipedia, Kalenjin languages (covers Keiyo as part of the cluster): \url{https://en.wikipedia.org/wiki/Kalenjin_languages}
\end{itemize}

\section{Religious and Community Resources}
\begin{itemize}
\item Joshua Project, Keiyo language page (population and Bible translation status, listed as started): \url{https://joshuaproject.net/languages/eyo}
\item Scripture Earth, Kalenjin (kln) scripture resources, a related macrolanguage, not Keiyo-specific: \url{https://www.scriptureearth.org/00eng.php?iso=kln}
\end{itemize}

\section{Digital and Audio Resources}
\begin{itemize}
\item Global Recordings Network, Keiyo (Geyo) audio-visual Bible lessons and evangelistic recordings: \url{https://globalrecordings.net/en/language/3019}
\end{itemize}

\section{Curriculum Development Context}
Keiyo is vigorous but not known to be taught in schools, and written materials appear scarce and are mostly religious. Kalenjin-level materials may partly serve as a base, but dialect differences need checking with Keiyo speakers.

\end{document}